\documentclass[10pt]{article}
\usepackage[margin=22mm]{geometry}
\usepackage{amsmath,amssymb,amsthm,booktabs,hyperref}
\newtheorem{theorem}{Theorem}
\newtheorem{lemma}[theorem]{Lemma}
\newtheorem{proposition}[theorem]{Proposition}
\newcommand{\GL}{\operatorname{GL}}
\newcommand{\GU}{\operatorname{GU}}
\newcommand{\Sp}{\operatorname{Sp}}
\newcommand{\PSL}{\operatorname{PSL}}
\newcommand{\PSU}{\operatorname{PSU}}
\newcommand{\PSp}{\operatorname{PSp}}
\newcommand{\Om}{\Omega}
\newcommand{\ii}{\mathrm i_2}
\newcommand{\F}{\mathbb F}
\setlength{\emergencystretch}{2em}
\title{Explicit involution classes for the classical part of L045}
\author{Private proof component: publication\_recovery}
\date{}
\begin{document}
\maketitle
\noindent This note proves the lower estimates needed in Section~4 of the
fixed L045 source, and corrects its special unitary degree-four estimate.
It does not establish the required maximal-overgroup or automorphism bounds,
or assert a complete recovery of the original paper. Throughout, $\ii(H)$
counts elements of order exactly two. The finite $35$ parameter certificates
are a separate, already distinctly reviewed component.

\begin{theorem}\label{main}
For every prime power $q$, the following strict lower bounds hold:
\[
\begin{array}{c|c|c}
G&\text{range}&\ii(G)>q^D/4\\\hline
\PSL_n(q),\ \PSU_n(q)&n\ge4&D=\lfloor n^2/2\rfloor\\
\PSp_{2m}(q)&m\ge2&D=m^2+m\\
\Om_{2m+1}(q)&m\ge3,\ q\text{ odd}&D=m^2+m\\
\operatorname{P}\Om_{2m}^{\epsilon}(q)&m\ge4,\ \epsilon=\pm1&D=m^2-1.
\end{array}
\]
The nonsimple group $\Sp_4(2)$ is allowed in this numerical statement;
it is not a positive simple-group obligation in L045.
\end{theorem}

\paragraph{Classical conventions and elementary counts.}
We use the natural finite-field matrix groups. In odd characteristic,
$\Om$ is the kernel of the spinor norm in the determinant-one orthogonal
group; in even characteristic it is the kernel of the Dickson invariant
$g\mapsto\operatorname{rank}(g-I)\bmod2$. These are the usual matrix
descriptions of the groups denoted $\Om$ in the source. A projective group
means the quotient by the scalar centre. The ordinary finite-field form
conventions, including the uniqueness of a nondegenerate hermitian form
in each dimension and the two signs of an even-dimensional quadratic form,
are retained elementary classical-group inputs. No maximal-subgroup or
involution-class classification is used below.

Put $P_j(q)=\prod_{i=1}^j(1-q^{-2i})$ and
$F_j(q)=\prod_{i=1}^j(1-(-q)^{-i})$, with empty products one. Counting
ordered bases, symplectic bases, orthogonal bases, and unitary bases gives
\begin{align*}
|\GL_j(q)|&=q^{j^2}\prod_{i=1}^j(1-q^{-i}),&
|\GU_j(q)|&=q^{j^2}F_j(q),\\
|\Sp_{2j}(q)|&=q^{2j^2+j}P_j(q),&
|O_{2j}^{\epsilon}(q)|&=2q^{2j^2-j}(1-\epsilon q^{-j})P_{j-1}(q),\\
|O_{2j+1}(q)|&=2q^{2j^2+j}P_j(q)&& (q\text{ odd}).
\end{align*}
We count whole ambient conjugacy orbits of matrices lying in the desired
special group. They need not be single conjugacy classes of that special
group. If $X^2$ is scalar, at most two scalar multiples of $X$ can lie in
its ambient class: $(zX)^2=X^2$ forces $z^2=1$. If the two eigenvalue
multiplicities are unequal, only $z=1$ is possible. Thus the divisions by
two below are valid lower bounds even when special-group classes split.

\begin{lemma}\label{products}
For $q\ge2$, $F_j(q)\le1+q^{-1}$. A consecutive interval of its factors
starting at $a$ has product at least $1-q^{-a}$, and has product greater
than one if $a$ is odd. Also
\[
\prod_{i\ge a}(1-q^{-2i})>1-\frac{q^{-2a}}{1-q^{-2}},\qquad
\prod_{\substack{i\ge1\\i\text{ odd}}}(1-q^{-i})
>1-\frac q{q^2-1}.
\]
In the last inequality we use $q\ge3$; the finite products used below obey
the same strict estimates.
\end{lemma}
\begin{proof}
Pair an even factor $1-q^{-2i}$ with its following odd factor to obtain a
product less than one; this proves the upper bound for $F_j$. Pair an odd
factor with its following even factor to obtain a product greater than
one; an interval beginning at an even index loses only its first factor.
The remaining inequalities follow by expanding a product of numbers
$1-a_i$, $0<a_i<1$, and bounding it below by $1-\sum a_i$.
\end{proof}

\paragraph{Linear and unitary groups in odd characteristic.}
For odd $n$ choose an even $k$ nearest $n/2$ and take the matrix with
$-1$-eigenspace of dimension $k$. For even $n=2m$ with even $m$, take
the balanced $m+m$ decomposition. All these matrices have determinant
one; use orthogonal hermitian summands in the unitary case. Their ambient
orbits have sizes $|\GL_n|/(|\GL_k||\GL_{n-k}|)$ or
$|\GU_n|/(|\GU_k||\GU_{n-k}|)$, respectively. The normalized linear
ratio exceeds one. For the unitary ratio, $k,n-k\ge2$, so Lemma~\ref{products}
gives
\[
\frac{F_n(q)}{F_k(q)F_{n-k}(q)}
>\frac{1-q^{-3}}{1+q^{-1}}\ge\frac{13}{18}.
\]
The orbit exponent is $2k(n-k)=\lfloor n^2/2\rfloor$. Division by at
most two still leaves a constant exceeding $1/4$.

If $n=2m$ and $m$ is odd, use a determinant-one matrix $X$ with
$X^2=-I$. In the linear case with $q\equiv1\pmod4$, take
$\operatorname{diag}(aI_m,-aI_m)$, $a^2=-1$; with $q\equiv3\pmod4$,
take multiplication by $a$ on $\F_{q^2}^m$. The centralizers are
$\GL_m(q)^2$ and $\GL_m(q^2)$. In the unitary case the two congruences
are reversed: $a$ has norm one if $q\equiv3\pmod4$, while for
$q\equiv1\pmod4$ the two eigenspaces are paired totally isotropic spaces,
and the centralizer is $\GL_m(q^2)$. In the latter cases the respective
normalized orbit ratios are
\[
\prod_{\substack{1\le i\le2m-1\\i\text{ odd}}}(1-q^{-i})>5/8,
\qquad
\prod_{\substack{1\le i\le2m-1\\i\text{ odd}}}(1+q^{-i})>1.
\]
The exponent is $2m^2$, and projectivization loses at most two.

\paragraph{Linear and unitary groups in characteristic two.}
Write $n=2k+l$, $l\in\{0,1\}$. A linear involution of Jordan shape
$(2^k,1^l)$ has centralizer
$q^{k(2n-3k)}\GL_k(q)\GL_l(q)$. This is obtained by writing commuting
endomorphisms along its image/kernel filtration: the invertible diagonal
blocks are one repeated $k$-block and one $l$-block. Its normalized orbit
ratio is
\[
\frac{\prod_{i=k+1}^n(1-q^{-i})}{\prod_{i=1}^l(1-q^{-i})}>3/4.
\]
For a hermitian space, put $X=I+N$. Unitarity and $N^2=0$ imply that
$W=\operatorname{im}N$ is totally isotropic and $N$ is self-adjoint.
Conversely, for each isotropic $k$-space $W$, the rank-$k$ self-adjoint
maps $V/W^{\perp}\longrightarrow W$ are precisely the nonsingular
hermitian matrices; their number is $|\GL_k(q^2)|/|\GU_k(q)|$.
The stabilizer of $W$ has order
$q^{k(2n-3k)}|\GL_k(q^2)||\GU_l(q)|$: the radical has $2kl+k^2$
parameters over $\F_q$, and the indicated two diagonal blocks are free.
Therefore the total number is
\[
\frac{|\GU_n(q)|}{q^{k(2n-3k)}|\GU_k(q)||\GU_l(q)|}.
\]
Its exponent is $2k(n-k)$ and its normalized ratio
$F_n/(F_kF_l)$ exceeds $(1-2^{-3})/(1+2^{-1})=7/12$.
All matrices just counted are unipotent, hence determinant one. In
characteristic two no nontrivial scalar multiple of a unipotent matrix
is unipotent. This proves both projective lower bounds.

\paragraph{Symplectic groups in odd characteristic.}
Use a symplectic complex structure $X^2=-I$. If $q\equiv1\pmod4$, its
two isotropic eigenspaces are paired and its centralizer is $\GL_m(q)$.
If $q\equiv3\pmod4$, view the space as a hermitian $m$-space over
$\F_{q^2}$ and obtain centralizer $\GU_m(q)$. These statements follow
directly by requiring a commuting map to preserve the pairing or the
hermitian form, respectively. The orbit exponent is $m^2+m$; the normalized
ratios are
\[
\prod_{i=1}^m(1+q^{-i})>1,
\qquad
\prod_{i=1}^m(1+(-q)^{-i})>5/8.
\]
At most two-to-one scalar fusion gives the result.

\paragraph{Symplectic groups in characteristic two.}
Fix an isotropic $k$-space $W$. The symplectic rank-$k$ square-zero maps
with image $W$ correspond to nonsingular symmetric $k\times k$ matrices.
Retain just the nonalternating ones. Every nonalternating nondegenerate
symmetric form over the perfect field $\F_q$ is equivalent to $I_k$.
One may diagonalize after splitting a nonsingular vector; a remaining
alternating plane together with one nonsingular line is equivalent to
three nonsingular lines. Its isometry group has order
\[
J_k(q)=\begin{cases}|\Sp_{k-1}(q)|,&k\text{ odd},\\
q^{k-1}|\Sp_{k-2}(q)|,&k\text{ even}.
\end{cases}
\]
Indeed the distinguished vector representing $v\mapsto\sqrt{(v,v)}$
is fixed. If it is nonsingular its perpendicular complement is symplectic;
if it is singular the quotient of its perpendicular space is symplectic
and the kernel, by the two remaining row/column equations, has $q^{k-1}$
elements. The isotropic-space stabilizer has Levi part
$\GL_k(q)\Sp_{2m-2k}(q)$ and radical order $q^u$, where
$u=(4mk-3k^2+k)/2$ (count the free off-diagonal block and its symmetric
correction). Consequently the number of these involutions is exactly
\begin{equation}\label{spcount}
C_{m,k}(q)=\frac{|\Sp_{2m}(q)|}{q^u|\Sp_{2m-2k}(q)|J_k(q)}.
\end{equation}
For $k=m$ this is $q^{m^2+m}$ times
$\prod_{i=\lceil m/2\rceil}^m(1-q^{-2i})>2/3$. The symplectic centre is
trivial in characteristic two.

\paragraph{Orthogonal groups in odd characteristic.}
For odd dimension $2m+1$, let $2a$ be the even member of $\{m,m+1\}$,
and $b=m-a$. Choose a nondegenerate $2a$-space with square discriminant;
such a space of one of the two signs exists with a nondegenerate
complement. Acting as $-I$ on it and $I$ on its complement has determinant
and spinor norm one. The full orthogonal orbit has size
\[
\frac{|O_{2m+1}(q)|}{|O_{2a}^{\delta}(q)||O_{2b+1}(q)|}.
\]
All its elements lie in $\Om$. The exponent is $m^2+m$, and its normalized
ratio exceeds $(1/2)(7/8)/(4/3)=21/64$.

For dimension $2m$, put $a=\lfloor m/2\rfloor$, $b=m-a$, choose again
a square-discriminant $2a$-space, and use the same sign involution. Its
orthogonal orbit has size
\[
\frac{|O_{2m}^{\epsilon}(q)|}
{|O_{2a}^{\delta}(q)||O_{2b}^{\epsilon\delta}(q)|}.
\]
The exponent is $4ab$. Apart from the leading factor $1/2$, its normalized
ratio is greater than
\[
\frac{(1-3^{-4})(7/8)}{(1+3^{-2})^2}=7/10.
\]
If $m$ is odd the two multiplicities differ, so the projective map is
injective on this orbit and $4ab=m^2-1$. If $m$ is even it loses at most
two and $4ab=m^2$; the extra factor $q\ge3$ more than pays this loss.

\paragraph{Orthogonal groups in characteristic two.}
Use \eqref{spcount} with even $k=m$ for even $m$, and $k=m-1$ for odd $m$.
In a hyperbolic symplectic basis the chosen involution sends $f_i$ to
$f_i+e_i$ for $1\le i\le k$ and fixes all other basis vectors. A quadratic
form $Q$ with this polar form is invariant precisely when $Q(e_i)=1$
for these $k$ indices. Thus exactly $q^{2m-k}$ forms are invariant.
Their Arf trace is
$\operatorname{Tr}(\sum Q(e_i)Q(f_i))$; since $k>0$, one free $Q(f_i)$
makes its two values equally frequent. Exactly $q^{2m-k}/2$ invariant
forms have each sign. There are $q^m(q^m+\epsilon)/2$ forms of sign
$\epsilon$ with the fixed polar form. Double counting pairs $(X,Q)$ gives
\begin{equation}\label{orthcount}
C_{m,k}(q)\frac{q^{m-k}}{q^m+\epsilon}
\end{equation}
involutions in each orthogonal group. Their Dickson invariant is $k\bmod2=0$,
so all belong to $\Om$; its scalar centre is trivial here. For even $m$,
\eqref{orthcount} is
\[
q^{m^2}\frac{\prod_{i=m/2}^m(1-q^{-2i})}{1+\epsilon q^{-m}}.
\]
For odd $m$ it is
\[
q^{m^2-1}\frac{\prod_{i=(m-1)/2}^m(1-q^{-2i})}
{(1-q^{-2})(1+\epsilon q^{-m})}.
\]
Both normalized ratios exceed $(11/12)/(17/16)=44/51$. This finishes
Theorem~\ref{main}.

\begin{proposition}\label{u4}
If $q\equiv3\pmod4$, then
\[
\ii(\PSU_4(q))=
\frac{q^4(q^2-q+1)(q^2+1)}2+
\begin{cases}0,&q\equiv3\pmod8,\\q^3(q-1)(q^2+1),&q\equiv7\pmod8.
\end{cases}
\]
In particular the source's special assertion $\ii(\PSU_4(q))>q^8/2$
is false for all such $q$.
\end{proposition}
\begin{proof}
A projective involution lifts to $X\in\operatorname{SU}_4(q)$ with
$X^2=zI$, $z\in\mu_4$. It is semisimple. For $z=1,-1$, determinant one
forces the noncentral multiplicities to be $2+2$. Both choices give the
same projective family; its size is
$|\GU_4|/(2|\GU_2|^2)$, the first summand displayed above.
For $z$ of order four, choose $a^2=z$, so $a$ has order eight. If
$q\equiv3\pmod8$, $a^{q+1}=-1$; its two eigenspaces must be paired
isotropic spaces of dimensions $2+2$. But the determinant is then
$a^4=-1$, impossible. If $q\equiv7\pmod8$, $a^{q+1}=1$; the eigenspaces
are nondegenerate and determinant one requires multiplicities $1+3$ or
$3+1$. Scalar multiplication cycles all four choices of $a$ and identifies
the two multiplicity orientations, leaving one projective family. Within
one $1+3$ ambient orbit the projective map is injective, so its size is
$|\GU_4|/(|\GU_1||\GU_3|)=q^3(q-1)(q^2+1)$.
There are no remaining possibilities for $z$ or for the eigenspaces.

For $q\equiv3\pmod8$, subtracting $q^8/2$ from the first summand gives
$-(q^4/2)(q(q-1)^2-1)<0$. For $q\equiv7\pmod8$, the difference is
$-(q^3/2)(q^4-4q^3+3q^2-3q+2)<0$, since $q\ge7$.
\end{proof}

\paragraph{Exact controls and attribution.}
The accompanying standard-library program evaluates all written class
formulas with rational arithmetic at $1,619$ parameter choices, verifies
the strict quarter bounds, and checks the correction at seven fields.
These finite checks are controls on the formulas, not premises for the
uniform inequalities. The ordinary character table $\mathrm U_4(3)$ gives
$2,835$ involutions, agreeing with Proposition~\ref{u4}; at $q=7$ the
formula gives $2,683,975<7^8/2$. A fresh clean table replay is retained with the component. The construction and calculations in this note
are new recovery work by publication\_recovery, responding to root's
request for explicit lower estimates. The fixed source is Leemans--Liebeck,
\emph{Chiral polyhedra and finite simple groups}, arXiv:1603.07713v3,
Lemma~2.4 and Section~4 (especially printed page~9). No claim of novelty or
priority is made for the classical formulas. The remaining existence
argument must separately pay every bad-generator and reverser bound.
\end{document}
